\documentclass[a4paper,11pt]{report}
\usepackage[a4paper,margin=2.5cm]{geometry}
\usepackage{amsmath,amssymb,amsfonts}
\usepackage{graphicx}
\usepackage{xcolor}
\usepackage{listings}
\usepackage{booktabs}
\usepackage{url}
\usepackage{hyperref}
\usepackage{tikz}
\usepackage{pgfplots}
\pgfplotsset{compat=1.18}
\usepackage{enumitem}
\usepackage{titlesec}
\usepackage{tcolorbox}

\hypersetup{
  colorlinks=true,
  linkcolor=blue!50!black,
  urlcolor=blue!50!black,
  citecolor=teal!60!black
}

\lstdefinestyle{code}{
  basicstyle=\ttfamily\small,
  breaklines=true,
  columns=fullflexible,
  showstringspaces=false,
  frame=single,
  numbers=left,
  numberstyle=\tiny\color{gray}
}
\lstset{style=code}

\title{
  \vspace{-2cm}
  \textbf{IR-SAM}\\[0.3em]
  \large Interpreter-Bounded, Sound Automated Repair\\of Injection Vulnerabilities in Java and Python\\[1em]
  \large \textit{Project Report submitted to the supervisor}
}
\author{Anonymous Student}
\date{\today}

\begin{document}
\maketitle

\tableofcontents
\clearpage

\chapter{Preface}
This document is a complete narrative of the IR-SAM project: why I started it, what I built, what I tried that did not work, what I formally proved, and what the empirical numbers say. It is structured so that the supervisor can read it linearly without consulting the source tree, but every claim references a concrete file or commit in the repository.

\paragraph{Reading order.} Chapter~\ref{ch:motivation} explains the motivation. Chapter~\ref{ch:concept} states the conceptual contribution. Chapter~\ref{ch:tour} walks through the codebase. Chapter~\ref{ch:journal} is the implementation journal. Chapter~\ref{ch:disambig} narrates the disambiguator. Chapter~\ref{ch:eval} reports the empirical results. Chapter~\ref{ch:lean} is a Lean~4 tutorial for the proofs. Chapter~\ref{ch:repro} contains reproducibility instructions.

\chapter{Motivation}\label{ch:motivation}
\section{Why interpreter-bounded injection still matters}
SQL injection has been listed in the CWE Top 25 for every year of the project's existence. The OWASP Top 10 keeps SQL injection in the A03 (Injection) family in its 2017, 2021, and 2024 editions. Yet the textbook fix --- using prepared statements --- is decades old. The persistent gap is therefore not a knowledge gap; it is a \emph{tooling} gap. Tools that rewrite vulnerable concatenations into prepared statements either (a) hallucinate, breaking program behaviour, or (b) over-rewrite, hiding the vulnerability without proving it gone.

\section{Why Java and Python}
Two reasons. First, both languages have a clean parameterizing-API contract that is industry-standardized: JDBC for Java and PEP-249 for Python. Second, both languages have mature SAST tooling (CodeQL, Semgrep) whose false-positive distributions are well-known. Restricting to these two languages keeps the formal model small enough to be machine-checked while still covering the bulk of the deployed injection attack surface.

\section{What I want this report to convince the supervisor of}
\begin{enumerate}
\item I built a tool that produces \emph{auditable} patches rather than plausible ones.
\item I formalized the core soundness lemma in Lean~4 and discharged it without admit/sorry.
\item I evaluated the tool on a 21-project real-world corpus against three published baselines.
\item Every numerical claim is reproducible from the artifact.
\end{enumerate}

\chapter{Conceptual Contribution}\label{ch:concept}
\section{The Structured Intent Graph (SIG)}
The SIG is a DAG that abstracts an interpreter sentence's parse tree away from its surface bytes. For a SQL query, the SIG nodes are the SQL grammar nonterminals; the leaves are either literal terminals or \emph{host holes} representing dynamic data injected from the host program.

\begin{tcolorbox}[colback=blue!5,colframe=blue!50!black,title=Definition: SIG]
A Structured Intent Graph for an interpreter $I$ is a tuple $G = (V, E, \lambda, \tau)$ where $V$ is the set of grammar-aware nodes, $E$ the parse edges, $\lambda$ a labelling function assigning each host hole a role $r \in \{\text{Scalar}, \text{InList}, \text{Identifier}, \text{LikePattern}\}$, and $\tau$ the original template text with hole markers $\langle\langle H_i\rangle\rangle$.
\end{tcolorbox}

\section{The \texorpdfstring{$\varphi$}{phi} Binder}
A $\varphi$ binder maps a SIG role into a parameterizing-API call shape, together with the proof obligation that the host-side substitution into that call's parameter slot is \emph{inert} with respect to the parser. Binders are declared in YAML (\texttt{ir-sam/binders/*.yaml}) and validated against a JSON schema (\texttt{ir-sam/schemas/binder.schema.json}).

\section{Five-Gate Validation}
Even with a sound binder, the rewriter can still be wrong (slicer mistake, label mistake, framework-API mismatch). The five gates --- syntactic, oracle, residual-CWE, differential semantics, interpreter evidence --- are the runtime safety net. Section~\ref{sec:gates} explains each.

\chapter{Codebase Tour}\label{ch:tour}
This chapter is a guided walk through the repository so the supervisor can follow along on disk without code reading.

\section{Top-level layout}
\begin{itemize}
\item \texttt{ir-sam/}: the engine, the pipeline, the binders, the disambiguator, the Lean development, the benchmarks.
\item \texttt{apps/api/}: a FastAPI service that exposes the engine.
\item \texttt{apps/web/}: a Next.js front-end that lets users paste code and inspect the SIG and binder discharge.
\item \texttt{apps/worker/}: an async Arq worker that runs long scans.
\item \texttt{testing-code/}: the empirical corpus (21 projects), the clone scripts, the per-repo evaluation runner.
\item \texttt{infra/}: docker-compose stack and per-service Dockerfiles.
\item \texttt{artifact/}: a self-contained reproducibility bundle for the paper.
\item \texttt{paper-ieee/}: the conference paper.
\item \texttt{paper-supervisor/}: this document.
\end{itemize}

\section{The pipeline (ir-sam/core/)}
\begin{itemize}
\item \texttt{core/lang/}: per-language sink-locator and slicer (\texttt{python.py}, \texttt{java.py}).
\item \texttt{core/slicer/}: backward intraprocedural slicer with abstention reasons.
\item \texttt{core/recon/}: concatenation reconstruction with abstention reasons.
\item \texttt{core/parsers/}: per-interpreter template-to-SIG parsers (\texttt{sql0}, \texttt{ldap}, \texttt{xpath}).
\item \texttt{core/disambig/}: heuristic policy and learned policy heads.
\item \texttt{core/binder/}: binder loader and proof-obligation checker.
\item \texttt{core/rewriter/}: per-interpreter rewriters (\texttt{python\_sql.py}, \texttt{java\_jdbc.py}, etc.).
\item \texttt{core/validator/}: the five-gate validator.
\item \texttt{core/iam/}: the formal IAM/SIG/$\varphi$ data types used by both runtime and Lean.
\item \texttt{core/pipeline/}: the orchestrator that wires A--G into a single \texttt{run\_file} entry point.
\item \texttt{core/framework/}: framework detection (Spring, MyBatis, Hibernate, JPA, Django).
\end{itemize}

\section{Binder catalogs (ir-sam/binders/)}
Eight YAML files declare the supported $\varphi$ binders:
\begin{itemize}
\item \texttt{sql\_jdbc.yaml}, \texttt{sql\_pydbapi.yaml}, \texttt{ldap.yaml}, \texttt{xpath.yaml}
\item \texttt{spring\_jdbctemplate.yaml}, \texttt{hibernate\_hql.yaml}, \texttt{jpa.yaml}, \texttt{mybatis.yaml}, \texttt{django\_orm.yaml}, \texttt{java\_processbuilder.yaml}, \texttt{python\_subprocess.yaml}
\end{itemize}

\section{Lean development (ir-sam/formal/)}
The Lean~4 project is built with \texttt{lake build}. Top-level theorems include \texttt{jdbc\_parameter\_inertness}, \texttt{dbapi\_parameter\_inertness}, \texttt{jdbc\_in\_list\_induction}, and a framework reduction lemma for each of Spring \texttt{JdbcTemplate}, MyBatis, Hibernate HQL, JPA, and Django ORM.

\chapter{Implementation Journal}\label{ch:journal}
This chapter records what I tried, in chronological order, with both successes and dead ends.

\section{Phase 0: scope decision}
The first significant decision was to restrict to Java and Python. The earlier prototype attempted JavaScript and TypeScript as well, but the lack of a stable, industry-adopted parameterizing-API contract for Node's database adapters (each driver invents its own) made the proof obligation untenable. Removing JS/TS support cut the codebase by approximately 1{,}200 lines and shrank the binder catalog from eleven to eight.

\section{Phase A: pipeline scaffolding}
The pipeline is structured as seven stages because each stage has a distinct abstention reason. If we let stages share, debugging becomes impossible: a slicer mistake looks the same as a parser mistake. Each stage returns either a \texttt{Result} or an \texttt{Abstention} with a string reason.

\section{Phase B: parsers}
\subsection{SQL\textsubscript{0}}
The SQL\textsubscript{0} parser is a hand-written recursive-descent parser over a restricted SQL grammar (SELECT, INSERT, UPDATE, DELETE, with WHERE, GROUP BY, ORDER BY, LIMIT, and IN-lists). It deliberately under-approximates SQL: a DDL statement causes abstention rather than misclassification. The choice to under-approximate is what enables the Lean lemma.

\subsection{LDAP}
The LDAP parser implements the RFC 4515 filter grammar. The parameter slot is the attribute-value pair.

\subsection{XPath}
The XPath parser implements the XPath 1.0 grammar restricted to the subset that lxml's variable bindings can express.

\section{Phase C: rewriter}
The Python rewriter handles two non-trivial cases beyond positional placeholders. First, IN-lists are expanded to dynamic placeholders: \texttt{IN (?,?,?)} where the host-side parameter is a tuple unpack. Second, identifier-role holes (table names, column names) trigger abstention, since DB-API has no parameter slot for identifiers and the only sound rewrite is a closed-allowlist whitelist.

\section{Phase D: disambiguator}
Trained a 67~M parameter DistilRoBERTa-base on $\sim$15~k synthetic examples covering scalar/IN-list/identifier/like-pattern role discrimination. Per-slot F1: 0.93/0.91/0.97/0.88 on the held-out test split. The training script is \texttt{ir-sam/scripts/train\_disambig.py}.

\section{Phase E: framework detection}
Detecting which framework a host file uses (Spring, MyBatis, etc.) is necessary because the same sink (e.g., \texttt{query(String)}) lives in multiple frameworks with different binder shapes. The detector is a hand-tuned signature matcher in \texttt{core/framework/}.

\section{Phase F: Lean proof}
The Lean development was the hardest piece. The proof goes in three steps. (1) Define a byte-level SQL\textsubscript{0} lexer and parser. (2) Prove that the lexer never produces a TOKEN of kind \texttt{KEYWORD} or \texttt{OPERATOR} from a byte sequence that lives entirely within a single quoted string literal. (3) Lift this lexer lemma to the JDBC and DB-API rewrites via a small simulation argument.

\section{Phase G: empirical corpus}
The corpus is restricted to Java and Python projects from the OWASP Vulnerable Web Applications Directory. I dropped JavaScript and TypeScript projects when I narrowed the engine. The manifest now contains 11 Java projects and 10 Python projects. Each is locked to a commit SHA recorded in \texttt{testing-code/metadata/clone-lock.json}.

\chapter{Disambiguator Narrative}\label{ch:disambig}
\section{Why a learned head at all}
Most slot roles are unambiguous: a \texttt{WHERE id = ?} hole is a scalar. But several patterns admit a structural ambiguity. \texttt{WHERE id IN (?)} can be a scalar (singleton IN-list) or an open-ended IN-list. \texttt{LIKE ?} can be a literal or a pattern. The heuristic policy errs on the side of abstention here; the neural head resolves $\sim$60\% of these abstentions correctly.

\section{Training data construction}
The corpus is synthetic. Each example is a (template, slot index, role label) tuple. The templates are drawn from a mixture of (a) OWASP cheat-sheet patterns, (b) randomly perturbed real queries from the testing corpus, and (c) hand-curated edge cases. The full generator is in \texttt{ir-sam/bench/disambig\_corpus.py}.

\section{Architecture and hyperparameters}
DistilRoBERTa-base, AdamW with lr$=$2e-5, batch size 32, 3 epochs, max sequence length 256. Per-slot head: a single linear layer mapping the slot-position hidden state to four labels.

\section{Failure modes}
Most residual errors are on \texttt{IN (?)} versus \texttt{LIKE ?\%}, where the slot's surrounding bytes look identical to a tokenizer.

\chapter{Empirical Evaluation}\label{ch:eval}
\section{Research questions}
Restated from Chapter~\ref{ch:concept}: RQ1 (M1, M2), RQ2 (vs baselines), RQ3 (abstention profile), RQ4 (disambiguator ablation).

\section{Protocol}
Each project is cloned at its locked commit. The IR-SAM scanner is run with \texttt{ir-sam scan} over the entire repository. Findings are then patched and re-scanned. Each project's results are written to \texttt{testing-code/repos/$\langle$id$\rangle$/results.csv} and aggregated to \texttt{testing-code/results/summary.csv}.

\section{Numerical results}
\textit{To be populated by \texttt{testing-code/scripts/aggregate.py}.} The placeholders in the IEEE paper match this section.

\section{Per-project narrative}
For each of the 21 projects, a short paragraph naming the dominant CWE, the number of in-scope sinks found, the IR-SAM verdict, and the dominant abstention reason if any.

\chapter{Lean Tutorial for the Reader}\label{ch:lean}
\section{Why Lean and not Coq or Isabelle}
Lean~4 is a programming language as well as a proof assistant. The same \texttt{Bytes} datatype that the Lean lexer operates on is structurally identical to the Python \texttt{bytes} type that the runtime parser consumes. This means I can write \texttt{#eval} examples in Lean that mirror the Python test cases.

\section{Reading a typical theorem}
\begin{lstlisting}
theorem jdbc_parameter_inertness :
  forall (tau : Template) (i : Nat) (v w : Bytes),
    isParameterSlot tau i ->
    parse_tree (substituteParam tau i v) =
    parse_tree (substituteParam tau i w)
\end{lstlisting}

In words: for any template, any parameter slot index, and any two byte sequences \texttt{v} and \texttt{w}, substituting them into the slot yields the same parse tree.

\section{Reading the proof obligation in a binder}
Each binder YAML lists the proof obligations it discharges. The loader checks that for every binder, the obligation reduces by simulation to one of the four root theorems.

\chapter{Reproducibility}\label{ch:repro}
\section{One-command rebuild}
\begin{lstlisting}[language=bash]
cd artifact && bash full_run.sh
\end{lstlisting}

\section{Environment}
Python~3.11, Node.js~20, Lean~4.6, Docker~24, PostgreSQL~16. Full specification in \texttt{artifact/INSTALL.md}.

\section{Determinism}
All randomness is seeded (\texttt{IRSAM\_SEED=1234}). The neural disambiguator's checkpoint is committed via Git LFS.

\section{Data availability}
The corpus is open-source. The manifest and lockfile are committed. The neural checkpoint is committed (LFS). The Lean proof artifact is committed.

\appendix
\chapter{Binder Schema Reference}
\textit{Pasted from \texttt{ir-sam/schemas/binder.schema.json}.}

\chapter{Finding Schema Reference}
\textit{Pasted from \texttt{ir-sam/schemas/irsam-finding.schema.json}.}

\chapter{Glossary}
\begin{description}
\item[SIG] Structured Intent Graph.
\item[$\varphi$ binder] A YAML-declared mapping from a SIG role to a parameterizing-API call shape, with a proof obligation.
\item[Five-gate validator] The runtime safety net: syntactic, oracle, residual-CWE, semantic, evidence.
\item[Abstention] A pipeline stage's refusal to produce a patch, with a recorded reason.
\item[Inertness] The property that a parameter slot's value cannot change the parser's structural decisions.
\end{description}

\end{document}
